% Erratum E3 for inclusion in main.tex.
% Suggested use: % Erratum E3 for inclusion in main.tex.
% Suggested use: % Erratum E3 for inclusion in main.tex.
% Suggested use: % Erratum E3 for inclusion in main.tex.
% Suggested use: \input{erratum_E3_distance_identity}

\subsection*{Erratum E3: the distance-function identity (9)--(10)}
\label{s:erratum:E3}

The identity (9) is false as stated because the quantity $E$ in
(10) is not the right distance quantity.  Namely,
for any $r>0$ the vector $x_1=0$ is admissible in the infimum defining $E$,
and therefore
\[
	\inf_{\norm{x_1}{1}\leq r}
	\norm{x-x_1}{0}^{1-\theta}\norm{x_1}{1}^{\theta}
	=0 .
\]
Thus $E=0$ for every $x$.  In general this cannot equal
$N_\theta^{-1}\norm{x}{\theta:1}$.  For instance, if $X=X_1=\IR$, both with
the usual norm, and $x=1$, then
\[
	d(r)=\max\{1-r,0\},
	\qquad
	D=\sup_{0<r<1}(1-r)^{1-\theta}r^\theta
	=(1-\theta)^{1-\theta}\theta^\theta>0,
\]
whereas the printed quantity $E$ is zero.

The correct statement is
\begin{equation}
	N_\theta^{-1}\norm{x}{\theta:1}
	=
	D,
	\qquad
	D:=\sup_{r>0}d(r)^{1-\theta}r^\theta .
	\label{e:erratum:E3:correct-D}
\end{equation}
Consequently, step c) in the proof of (9) should be deleted.  The
passage from an infimum over $\norm{x_1}{1}\leq r$ to an infimum over
$\norm{x_1}{1}=r$ is not valid for the printed expression $E$.

We include a proof of the corrected identity \eqref{e:erratum:E3:correct-D}.
Put
\[
	F:=\norm{x}{\theta:1},
	\qquad
	h(r):=d(r)^2 .
\]
From the definitions of $K_t$ and $d$, one has
\begin{equation}
	K_t(x)^2
	=
	\inf_{r\geq0}\{h(r)+t^2r^2\} .
	\label{e:erratum:E3:K-distance}
\end{equation}
Indeed, the inequality ``$\geq$'' follows by taking
$r=\norm{x_1}{1}$ in the definition of $K_t$, while the reverse inequality
follows by approximating the infimum in the definition of $d(r)$.

First suppose that $D<\infty$.  Then, for every $r>0$,
\[
	d(r)\leq D^{1/(1-\theta)} r^{-\theta/(1-\theta)} .
\]
Using this in \eqref{e:erratum:E3:K-distance} and minimizing over $r>0$ gives
\[
	K_t(x)^2
	\leq
	\inf_{r>0}
	\left\{
	D^{2/(1-\theta)}r^{-2\theta/(1-\theta)}+t^2r^2
	\right\}
	=
	N_\theta^2D^2t^{2\theta} .
\]
Hence
\[
	F=\sup_{t>0}t^{-\theta}K_t(x)\leq N_\theta D .
\]
If $D=\infty$, this inequality is void.

It remains to prove the reverse estimate.  If
$\lim_{r\to\infty}d(r)>0$, then both $D$ and $F$ are infinite: the former is
immediate from the definition of $D$, and the latter follows from
\eqref{e:erratum:E3:K-distance} by letting $t\downarrow0$.  We may therefore
assume that $d(r)\to0$ as $r\to\infty$.

Let $r>0$ be such that $d(r)>0$.  Since $d$ is nonnegative, convex, and
nonincreasing, the function $h=d^2$ is convex and nonincreasing.  Choose a
subgradient $p\in\partial h(r)$.  Since $h$ is nonincreasing, every
subgradient is nonpositive.  Moreover no subgradient can be zero: if
$0\in\partial h(r)$, convexity would imply $h(s)\geq h(r)>0$ for all
$s\geq0$, contradicting $h(s)\to0$ as $s\to\infty$.  Hence $p<0$.
Define
\[
	t^2:=-\frac{p}{2r}>0 .
\]
The supporting-line inequality for $h$ gives, for all $s\geq0$,
\[
	h(s)\geq h(r)+p(s-r).
\]
Consequently,
\[
	h(s)+t^2s^2
	\geq
	h(r)+p(s-r)+t^2s^2
	=
	h(r)+t^2r^2+t^2(s-r)^2
	\geq
	h(r)+t^2r^2 .
\]
Thus $r$ minimizes $s\mapsto h(s)+t^2s^2$, and
\eqref{e:erratum:E3:K-distance} gives
\[
	K_t(x)^2=d(r)^2+t^2r^2 .
\]
With $u:=tr/d(r)>0$, it follows that
\[
	t^{-\theta}K_t(x)
	=
	d(r)^{1-\theta}r^\theta
	u^{-\theta}(1+u^2)^{1/2} .
\]
By the definition of $N_\theta$ in equation (4),
\[
	\inf_{u>0}u^{-\theta}(1+u^2)^{1/2}=N_\theta .
\]
Therefore
\[
	F\geq t^{-\theta}K_t(x)
	\geq
	N_\theta d(r)^{1-\theta}r^\theta .
\]
Taking the supremum over all $r>0$ with $d(r)>0$ yields
\[
	F\geq N_\theta D .
\]
Together with the first inequality, this proves
\eqref{e:erratum:E3:correct-D}.


\subsection*{Erratum E3: the distance-function identity (9)--(10)}
\label{s:erratum:E3}

The identity (9) is false as stated because the quantity $E$ in
(10) is not the right distance quantity.  Namely,
for any $r>0$ the vector $x_1=0$ is admissible in the infimum defining $E$,
and therefore
\[
	\inf_{\norm{x_1}{1}\leq r}
	\norm{x-x_1}{0}^{1-\theta}\norm{x_1}{1}^{\theta}
	=0 .
\]
Thus $E=0$ for every $x$.  In general this cannot equal
$N_\theta^{-1}\norm{x}{\theta:1}$.  For instance, if $X=X_1=\IR$, both with
the usual norm, and $x=1$, then
\[
	d(r)=\max\{1-r,0\},
	\qquad
	D=\sup_{0<r<1}(1-r)^{1-\theta}r^\theta
	=(1-\theta)^{1-\theta}\theta^\theta>0,
\]
whereas the printed quantity $E$ is zero.

The correct statement is
\begin{equation}
	N_\theta^{-1}\norm{x}{\theta:1}
	=
	D,
	\qquad
	D:=\sup_{r>0}d(r)^{1-\theta}r^\theta .
	\label{e:erratum:E3:correct-D}
\end{equation}
Consequently, step c) in the proof of (9) should be deleted.  The
passage from an infimum over $\norm{x_1}{1}\leq r$ to an infimum over
$\norm{x_1}{1}=r$ is not valid for the printed expression $E$.

We include a proof of the corrected identity \eqref{e:erratum:E3:correct-D}.
Put
\[
	F:=\norm{x}{\theta:1},
	\qquad
	h(r):=d(r)^2 .
\]
From the definitions of $K_t$ and $d$, one has
\begin{equation}
	K_t(x)^2
	=
	\inf_{r\geq0}\{h(r)+t^2r^2\} .
	\label{e:erratum:E3:K-distance}
\end{equation}
Indeed, the inequality ``$\geq$'' follows by taking
$r=\norm{x_1}{1}$ in the definition of $K_t$, while the reverse inequality
follows by approximating the infimum in the definition of $d(r)$.

First suppose that $D<\infty$.  Then, for every $r>0$,
\[
	d(r)\leq D^{1/(1-\theta)} r^{-\theta/(1-\theta)} .
\]
Using this in \eqref{e:erratum:E3:K-distance} and minimizing over $r>0$ gives
\[
	K_t(x)^2
	\leq
	\inf_{r>0}
	\left\{
	D^{2/(1-\theta)}r^{-2\theta/(1-\theta)}+t^2r^2
	\right\}
	=
	N_\theta^2D^2t^{2\theta} .
\]
Hence
\[
	F=\sup_{t>0}t^{-\theta}K_t(x)\leq N_\theta D .
\]
If $D=\infty$, this inequality is void.

It remains to prove the reverse estimate.  If
$\lim_{r\to\infty}d(r)>0$, then both $D$ and $F$ are infinite: the former is
immediate from the definition of $D$, and the latter follows from
\eqref{e:erratum:E3:K-distance} by letting $t\downarrow0$.  We may therefore
assume that $d(r)\to0$ as $r\to\infty$.

Let $r>0$ be such that $d(r)>0$.  Since $d$ is nonnegative, convex, and
nonincreasing, the function $h=d^2$ is convex and nonincreasing.  Choose a
subgradient $p\in\partial h(r)$.  Since $h$ is nonincreasing, every
subgradient is nonpositive.  Moreover no subgradient can be zero: if
$0\in\partial h(r)$, convexity would imply $h(s)\geq h(r)>0$ for all
$s\geq0$, contradicting $h(s)\to0$ as $s\to\infty$.  Hence $p<0$.
Define
\[
	t^2:=-\frac{p}{2r}>0 .
\]
The supporting-line inequality for $h$ gives, for all $s\geq0$,
\[
	h(s)\geq h(r)+p(s-r).
\]
Consequently,
\[
	h(s)+t^2s^2
	\geq
	h(r)+p(s-r)+t^2s^2
	=
	h(r)+t^2r^2+t^2(s-r)^2
	\geq
	h(r)+t^2r^2 .
\]
Thus $r$ minimizes $s\mapsto h(s)+t^2s^2$, and
\eqref{e:erratum:E3:K-distance} gives
\[
	K_t(x)^2=d(r)^2+t^2r^2 .
\]
With $u:=tr/d(r)>0$, it follows that
\[
	t^{-\theta}K_t(x)
	=
	d(r)^{1-\theta}r^\theta
	u^{-\theta}(1+u^2)^{1/2} .
\]
By the definition of $N_\theta$ in equation (4),
\[
	\inf_{u>0}u^{-\theta}(1+u^2)^{1/2}=N_\theta .
\]
Therefore
\[
	F\geq t^{-\theta}K_t(x)
	\geq
	N_\theta d(r)^{1-\theta}r^\theta .
\]
Taking the supremum over all $r>0$ with $d(r)>0$ yields
\[
	F\geq N_\theta D .
\]
Together with the first inequality, this proves
\eqref{e:erratum:E3:correct-D}.


\subsection*{Erratum E3: the distance-function identity (9)--(10)}
\label{s:erratum:E3}

The identity (9) is false as stated because the quantity $E$ in
(10) is not the right distance quantity.  Namely,
for any $r>0$ the vector $x_1=0$ is admissible in the infimum defining $E$,
and therefore
\[
	\inf_{\norm{x_1}{1}\leq r}
	\norm{x-x_1}{0}^{1-\theta}\norm{x_1}{1}^{\theta}
	=0 .
\]
Thus $E=0$ for every $x$.  In general this cannot equal
$N_\theta^{-1}\norm{x}{\theta:1}$.  For instance, if $X=X_1=\IR$, both with
the usual norm, and $x=1$, then
\[
	d(r)=\max\{1-r,0\},
	\qquad
	D=\sup_{0<r<1}(1-r)^{1-\theta}r^\theta
	=(1-\theta)^{1-\theta}\theta^\theta>0,
\]
whereas the printed quantity $E$ is zero.

The correct statement is
\begin{equation}
	N_\theta^{-1}\norm{x}{\theta:1}
	=
	D,
	\qquad
	D:=\sup_{r>0}d(r)^{1-\theta}r^\theta .
	\label{e:erratum:E3:correct-D}
\end{equation}
Consequently, step c) in the proof of (9) should be deleted.  The
passage from an infimum over $\norm{x_1}{1}\leq r$ to an infimum over
$\norm{x_1}{1}=r$ is not valid for the printed expression $E$.

We include a proof of the corrected identity \eqref{e:erratum:E3:correct-D}.
Put
\[
	F:=\norm{x}{\theta:1},
	\qquad
	h(r):=d(r)^2 .
\]
From the definitions of $K_t$ and $d$, one has
\begin{equation}
	K_t(x)^2
	=
	\inf_{r\geq0}\{h(r)+t^2r^2\} .
	\label{e:erratum:E3:K-distance}
\end{equation}
Indeed, the inequality ``$\geq$'' follows by taking
$r=\norm{x_1}{1}$ in the definition of $K_t$, while the reverse inequality
follows by approximating the infimum in the definition of $d(r)$.

First suppose that $D<\infty$.  Then, for every $r>0$,
\[
	d(r)\leq D^{1/(1-\theta)} r^{-\theta/(1-\theta)} .
\]
Using this in \eqref{e:erratum:E3:K-distance} and minimizing over $r>0$ gives
\[
	K_t(x)^2
	\leq
	\inf_{r>0}
	\left\{
	D^{2/(1-\theta)}r^{-2\theta/(1-\theta)}+t^2r^2
	\right\}
	=
	N_\theta^2D^2t^{2\theta} .
\]
Hence
\[
	F=\sup_{t>0}t^{-\theta}K_t(x)\leq N_\theta D .
\]
If $D=\infty$, this inequality is void.

It remains to prove the reverse estimate.  If
$\lim_{r\to\infty}d(r)>0$, then both $D$ and $F$ are infinite: the former is
immediate from the definition of $D$, and the latter follows from
\eqref{e:erratum:E3:K-distance} by letting $t\downarrow0$.  We may therefore
assume that $d(r)\to0$ as $r\to\infty$.

Let $r>0$ be such that $d(r)>0$.  Since $d$ is nonnegative, convex, and
nonincreasing, the function $h=d^2$ is convex and nonincreasing.  Choose a
subgradient $p\in\partial h(r)$.  Since $h$ is nonincreasing, every
subgradient is nonpositive.  Moreover no subgradient can be zero: if
$0\in\partial h(r)$, convexity would imply $h(s)\geq h(r)>0$ for all
$s\geq0$, contradicting $h(s)\to0$ as $s\to\infty$.  Hence $p<0$.
Define
\[
	t^2:=-\frac{p}{2r}>0 .
\]
The supporting-line inequality for $h$ gives, for all $s\geq0$,
\[
	h(s)\geq h(r)+p(s-r).
\]
Consequently,
\[
	h(s)+t^2s^2
	\geq
	h(r)+p(s-r)+t^2s^2
	=
	h(r)+t^2r^2+t^2(s-r)^2
	\geq
	h(r)+t^2r^2 .
\]
Thus $r$ minimizes $s\mapsto h(s)+t^2s^2$, and
\eqref{e:erratum:E3:K-distance} gives
\[
	K_t(x)^2=d(r)^2+t^2r^2 .
\]
With $u:=tr/d(r)>0$, it follows that
\[
	t^{-\theta}K_t(x)
	=
	d(r)^{1-\theta}r^\theta
	u^{-\theta}(1+u^2)^{1/2} .
\]
By the definition of $N_\theta$ in equation (4),
\[
	\inf_{u>0}u^{-\theta}(1+u^2)^{1/2}=N_\theta .
\]
Therefore
\[
	F\geq t^{-\theta}K_t(x)
	\geq
	N_\theta d(r)^{1-\theta}r^\theta .
\]
Taking the supremum over all $r>0$ with $d(r)>0$ yields
\[
	F\geq N_\theta D .
\]
Together with the first inequality, this proves
\eqref{e:erratum:E3:correct-D}.


\subsection*{Erratum E3: the distance-function identity (9)--(10)}
\label{s:erratum:E3}

The identity (9) is false as stated because the quantity $E$ in
(10) is not the right distance quantity.  Namely,
for any $r>0$ the vector $x_1=0$ is admissible in the infimum defining $E$,
and therefore
\[
	\inf_{\norm{x_1}{1}\leq r}
	\norm{x-x_1}{0}^{1-\theta}\norm{x_1}{1}^{\theta}
	=0 .
\]
Thus $E=0$ for every $x$.  In general this cannot equal
$N_\theta^{-1}\norm{x}{\theta:1}$.  For instance, if $X=X_1=\IR$, both with
the usual norm, and $x=1$, then
\[
	d(r)=\max\{1-r,0\},
	\qquad
	D=\sup_{0<r<1}(1-r)^{1-\theta}r^\theta
	=(1-\theta)^{1-\theta}\theta^\theta>0,
\]
whereas the printed quantity $E$ is zero.

The correct statement is
\begin{equation}
	N_\theta^{-1}\norm{x}{\theta:1}
	=
	D,
	\qquad
	D:=\sup_{r>0}d(r)^{1-\theta}r^\theta .
	\label{e:erratum:E3:correct-D}
\end{equation}
Consequently, step c) in the proof of (9) should be deleted.  The
passage from an infimum over $\norm{x_1}{1}\leq r$ to an infimum over
$\norm{x_1}{1}=r$ is not valid for the printed expression $E$.

We include a proof of the corrected identity \eqref{e:erratum:E3:correct-D}.
Put
\[
	F:=\norm{x}{\theta:1},
	\qquad
	h(r):=d(r)^2 .
\]
From the definitions of $K_t$ and $d$, one has
\begin{equation}
	K_t(x)^2
	=
	\inf_{r\geq0}\{h(r)+t^2r^2\} .
	\label{e:erratum:E3:K-distance}
\end{equation}
Indeed, the inequality ``$\geq$'' follows by taking
$r=\norm{x_1}{1}$ in the definition of $K_t$, while the reverse inequality
follows by approximating the infimum in the definition of $d(r)$.

First suppose that $D<\infty$.  Then, for every $r>0$,
\[
	d(r)\leq D^{1/(1-\theta)} r^{-\theta/(1-\theta)} .
\]
Using this in \eqref{e:erratum:E3:K-distance} and minimizing over $r>0$ gives
\[
	K_t(x)^2
	\leq
	\inf_{r>0}
	\left\{
	D^{2/(1-\theta)}r^{-2\theta/(1-\theta)}+t^2r^2
	\right\}
	=
	N_\theta^2D^2t^{2\theta} .
\]
Hence
\[
	F=\sup_{t>0}t^{-\theta}K_t(x)\leq N_\theta D .
\]
If $D=\infty$, this inequality is void.

It remains to prove the reverse estimate.  If
$\lim_{r\to\infty}d(r)>0$, then both $D$ and $F$ are infinite: the former is
immediate from the definition of $D$, and the latter follows from
\eqref{e:erratum:E3:K-distance} by letting $t\downarrow0$.  We may therefore
assume that $d(r)\to0$ as $r\to\infty$.

Let $r>0$ be such that $d(r)>0$.  Since $d$ is nonnegative, convex, and
nonincreasing, the function $h=d^2$ is convex and nonincreasing.  Choose a
subgradient $p\in\partial h(r)$.  Since $h$ is nonincreasing, every
subgradient is nonpositive.  Moreover no subgradient can be zero: if
$0\in\partial h(r)$, convexity would imply $h(s)\geq h(r)>0$ for all
$s\geq0$, contradicting $h(s)\to0$ as $s\to\infty$.  Hence $p<0$.
Define
\[
	t^2:=-\frac{p}{2r}>0 .
\]
The supporting-line inequality for $h$ gives, for all $s\geq0$,
\[
	h(s)\geq h(r)+p(s-r).
\]
Consequently,
\[
	h(s)+t^2s^2
	\geq
	h(r)+p(s-r)+t^2s^2
	=
	h(r)+t^2r^2+t^2(s-r)^2
	\geq
	h(r)+t^2r^2 .
\]
Thus $r$ minimizes $s\mapsto h(s)+t^2s^2$, and
\eqref{e:erratum:E3:K-distance} gives
\[
	K_t(x)^2=d(r)^2+t^2r^2 .
\]
With $u:=tr/d(r)>0$, it follows that
\[
	t^{-\theta}K_t(x)
	=
	d(r)^{1-\theta}r^\theta
	u^{-\theta}(1+u^2)^{1/2} .
\]
By the definition of $N_\theta$ in equation (4),
\[
	\inf_{u>0}u^{-\theta}(1+u^2)^{1/2}=N_\theta .
\]
Therefore
\[
	F\geq t^{-\theta}K_t(x)
	\geq
	N_\theta d(r)^{1-\theta}r^\theta .
\]
Taking the supremum over all $r>0$ with $d(r)>0$ yields
\[
	F\geq N_\theta D .
\]
Together with the first inequality, this proves
\eqref{e:erratum:E3:correct-D}.
